\chapter{Introduction}
\label{chap:introduction}
\section{Motivation}
Explaining and understanding topics has played a significant role in society throughout history.
This has remained a substantial challenge to the present day.
Throughout history, as technology has evolved, so too have methods of learning.



In recent years, the introduction of Large Language Models (LLMs)\cite{brown2020languagemodelsfewshotlearners},
such as ChatGPT, has fundamentally transformed educational landscapes by providing students with an immediately
accessible platform for knowledge acquisition and conceptual exploration.

However, the inherently "one-size-fits-all" architecture of these general-purpose models introduces significant
pedagogical risks in real-world academic environments.
Learner profiles, prior knowledge, and cognitive baselines are highly individualized, yet LLMs typically deliver
standardized outputs that fail to adapt to diverse learning styles, cultural contexts, or varying levels
of academic preparation.
Furthermore, these models lack the necessary adaptability to align with specific curricula, often overlooking
niche subjects, specialized terminology, or institution-specific course materials that fall outside their pretraining
scope or contextual limits.
This limitation is particularly critical given the high-stakes nature of education, where academic assessments and
grades can profoundly shape an individual’s career trajectory and long-term opportunities. Moreover,
the well-documented propensity of LLMs to generate plausible but factually incorrect information ("hallucinations")
is fundamentally unacceptable.
% Consequently, addressing these gaps through pedagogical alignment, domain-specific fine-tuning, and robust verification mechanisms is essential for the responsible, equitable, and effective integration of LLMs into modern academia.
\section{Problems}
In the previous paragraph, the core limitations of using general-purpose LLMs were mentioned;
these limitations are now examined in detail.
\begin{itemize}
    \item Hallucinations are a significant problem and challenge in using LLMs. These models may produce
          responses that appear coherent and rational but contain inaccurate information. Therefore, using
          them in education can pose a risk for students.
    \item Each individual possesses a unique scientific background. Some have cultivated greater
          expertise in scientific practices, while others have specialized more deeply in theoretical
          domains. One person may be familiar with numerous concepts that another has never
          encountered. Collectively, these variations underscore the necessity for language
          models to adapt their responses and generate personalized answers tailored to
          the specific needs and prior knowledge of each user.
          To illustrate, consider a student majoring in literature who encounters a question
          concerning artificial intelligence agents. In this scenario, the language model should
          provide clear explanations of foundational concepts related to computer science,
          artificial intelligence, and associated topics, thereby bridging the knowledge gap
          for this learner. By contrast, the same query posed by an undergraduate student in
          computer science would require a markedly different level of depth and technical detail.
          Moreover, such differences need not be confined to interdisciplinary contexts; they can be considerably more nuanced. Two students within the same academic discipline may still exhibit substantially divergent scientific backgrounds, owing to factors such as varying levels of attention and engagement with specific courses—one having studied a particular class more diligently while the other has not.
    \item Each class possesses its own unique and distinctive character. Even the same course
          taught by the same instructor in two different academic years can exhibit substantial
          differences in its details.
          Students frequently ask language models questions about specific class materials, yet
          the model remains entirely unaware of the particular content covered in that class.
          As a result, the language model tends to generate a generic, uniform response for all
          students. This represents a significant limitation, because the instructor may have
          explored a topic in considerably greater depth than the model’s explanation,
          or may have employed specialized terminology, notations, or conceptual frameworks
          unique to that course. Such discrepancies can lead students to receive incomplete or
          inaccurate information, potentially creating obstacles in their academic progress
          and future professional endeavors.
\end{itemize}
\section{Resulting Issues}
These problems can lead to several critical issues in educational applications of
general-purpose LLMs. The following consequences undermine the reliability and effectiveness
of such systems in supporting student learning:

\begin{itemize}
    \item Answers may include inaccuracies, as models draw upon broad but imperfect knowledge
          bases without verification against authoritative sources or course-specific materials.
    \item Some explanations may be irrelevant or missing key elements, failing to address
          the precise conceptual gaps or learning objectives relevant to a given learner.
    \item Answers may ignore details the instructor has already explained in class, resulting
          in misalignment between the model's output and the established curriculum narrative.
    \item Responses can sometimes contain hallucinations—plausible yet factually incorrect
          information—which poses a particular risk in high-stakes educational contexts where
          students may internalize erroneous concepts without critical scrutiny.
    \item The model fails to identify and explain prerequisite knowledge and fill background gaps.
          It does not determine the prerequisite knowledge gaps of students and
          attempts to explain general prerequisites, which are often inaccurate for any given
          individual.
\end{itemize}

These limitations collectively highlight the inadequacy of generic LLMs for nuanced pedagogical support. Without mechanisms for personalization, curricular alignment, and factual grounding, their deployment risks not only diminishing learning outcomes but also eroding trust in AI-assisted education. Addressing these challenges necessitates the development of more sophisticated, context-aware AI systems capable of adapting to individual learners, specific course content, and institutional requirements.

\section{Research Questions}

The primary objective of this thesis is to develop an intelligent tutoring system that enhances student learning while minimizing the risk of providing incorrect or misleading information. To achieve this, the research addresses the following key questions:

\begin{itemize}
    \item \textbf{RQ1} \label{RQ1}: How can the system dynamically adapt its responses to align with a student's current knowledge level and learning progress?

    \item \textbf{RQ2} \label{RQ2}: How can course-specific content and pedagogical materials be effectively integrated to ensure domain-relevant and accurate responses?

    \item \textbf{RQ3} \label{RQ3}: What techniques and architectural strategies can be employed to effectively reduce hallucinations in large language model-based tutoring systems?

    \item \textbf{RQ4} \label{RQ4}: How can the system generate more accurate, context-aware, and educationally valuable answers that support meaningful learning?
\end{itemize}

\section{Contribution}

The main contributions of this thesis are as follows:

\begin{enumerate}
    \item A Retrieval-Augmented Generation (RAG) pipeline tailored for educational applications is designed and implemented.
    \item Various embedding models are experimented with, and the impact of embedding dimensions on retrieval quality is analyzed.
    \item A language model is fine-tuned for answering student questions in a domain-specific context.
    \item TripartiteRAG\footnote{pri}, a novel approach that combines the above techniques, is proposed.
    \item A dataset of student questions with their corresponding answers is collected.
    \item A handcrafted dataset for answer and prompt is created.
\end{enumerate}

\section{Overview}
